% 第 13 题：细管测重力加速度装置（TikZ）
% 左：开口向上竖直放置，液柱高 h 封住空气柱 L1（底部封闭）；
% 右：旋转至水平，空气柱 L2，液柱仍高 h（右端开口）。
\begin{tikzpicture}[font=\small, scale=0.95]
  % ---------- 左：竖直细管 ----------
  % 管壁（开口向上，底部封闭），管宽 1，高 5.4
  \draw[thick] (0,0) -- (0,5.4);
  \draw[thick] (1,0) -- (1,5.4);
  \draw[thick] (0,0) -- (1,0);
  % 液柱（上部，高 h=1.6，打剖面线）
  \fill[pattern=north east lines] (0.02,3.8) rectangle (0.98,5.38);
  % 空气柱 L1（下部）
  \node[rotate=90] at (1.5,1.7) {$L_{1}$};
  \node[rotate=90] at (1.5,4.6) {$h$};
  \draw[<->] (2.0,0.1) -- (2.0,3.7);
  \draw[<->] (2.0,3.9) -- (2.0,5.3);
  % ---------- 右：水平细管 ----------
  % 管壁（左端封闭，右端开口），管高 1，长 7.6
  \draw[thick] (5.2,0) -- (12.8,0);
  \draw[thick] (5.2,1) -- (12.8,1);
  \draw[thick] (5.2,0) -- (5.2,1);
  % 液柱（右端，高 h=1.6 对应水平长度 1.6）
  \fill[pattern=north east lines] (10.6,0.02) rectangle (12.78,0.98);
  % 空气柱 L2（左端）
  \draw[<->] (5.5,-0.5) -- (10.3,-0.5);
  \node[below] at (7.9,-0.5) {$L_{2}$};
  \draw[<->] (10.9,-0.5) -- (12.5,-0.5);
  \node[below] at (11.7,-0.5) {$h$};
\end{tikzpicture}
